\ELhold

\ELsection{سرفصل‌ها}{سرفصل‌ها}

\begin{ELunit}

\ELfigure{m02-course-outline}{}

\end{ELunit}

\ELhold

\ELsection{میدان‌های الکتریکی ساکن}{میدان‌های الکتریکی ساکن}

\begin{ELunit}

\begin{itemize}
\tightlist
\item
  در الکتریسیته ساکن بارهای الکتریکی در حال سکون هستند و میدان‌های الکتریکی با زمان تغییر نمی‌کنند.

  \begin{itemize}
  \tightlist
  \item
    در این حالت میدان‌های مغناطیسی وجود ندارند.
  \end{itemize}
\item
  توسعه الکتریسیته ساکن در فیزیک مقدماتی عموماً با قانون تجربی کولمب در مورد نیروی بین دو بار آغاز می‌شود.
\item
  در این درس به جای دنبال کردن پیشرفت تاریخی الکتریسیته ساکن، موضوع را با فرض دیورژانس و کرل شدت میدان الکتریکی در فضای آزاد معرفی می‌نماییم.
\end{itemize}

\ELfigure{m02-map}{}

\begin{itemize}
\tightlist
\item
  فرضیات اصلی الکتریسیته ساکن که بر اساس آن‌ها می‌توان دیگر روابط و قوانین این حوزه را استخراج کرد کدامند؟
\item
  چگونه با استفاده از فرضیات فوق قوانین کولمب و گوس استخراج می‌شوند؟
\item
  میدان الکتریکی ناشی از یک بار نقطه‌ای، تعدادی بار نقطه‌ای گسسته و یک توزیع بار پیوسته چگونه محاسبه می‌شود؟
\end{itemize}

\end{ELunit}

\ELhold

\ELsection{فرضیات اصلی الکتریسیته ساکن در فضای آزاد}{فرضیات اصلی الکتریسیته ساکن در فضای آزاد}

\begin{ELjoin}

\begin{itemize}
\tightlist
\item
  \textbf{شدت میدان الکتریکی}

  \begin{itemize}
  \tightlist
  \item
    نیروی وارد بر یک بار کوچک آزمون هنگامیکه در میدان قرار داده می‌شود، تقسیم بر بار. \ELim{\ELto} \ELim{\mathrm{N/C}} (نیوتن بر کولمب) معادل \ELim{\mathrm{V/m}} (ولت بر متر)
  \end{itemize}
\end{itemize}

\begin{ELdefinition}{}

\ELdisp{\vect{E}\ELbr =\lim_{q\to0}\frac{\vect{F}}{q}\qquad\ELbr (\mathrm{V/m})}

\end{ELdefinition}

\end{ELjoin}

\begin{ELjoin}

\begin{itemize}
\tightlist
\item
  شدت میدان الکتریکی متناسب و هم‌جهت با نیرو است.
\item
  دو فرض اساسی الکتریسیته ساکن در فضای آزاد
\end{itemize}

\begin{ELimportant}{}

\ELdisp{\nabla\cdot\vect{E}\ELbr =\frac{\rho}{\epsilon_0}\qquad\ELbr \qquad\ELbr  \nabla\times\vect{E}\ELbr =0}

\end{ELimportant}

\end{ELjoin}

\begin{ELunit}

\begin{itemize}
\item
  \ELim{\rho}: چگالی حجمی بار \ELim{(\mathrm{C/m^3})}؛ \ELim{\epsilon_0}: گذردهی فضای آزاد \ELim{(\mathrm{F/m})}
\item
  \ELim{\vect{E}} سلونوئیدی نیست مگر اینکه \ELim{\rho=0}
\item
  \ELim{\vect{E}} غیرگردشی است
\item
  شکل انتگرالی معادله اول \ELim{\ELto} انتگرال‌گیری روی حجم دلخواه \ELim{V}
\end{itemize}

\ELdisp{\int_V\nabla\cdot\vect{E}\,dv\ELbr =\frac{1}{\epsilon_0}\int_V\rho\,dv}

\end{ELunit}

\begin{ELjoin}

\begin{itemize}
\tightlist
\item
  \ELim{\int_V\rho\,dv}: کل بار موجود در حجم \ELim{V} احاطه شده توسط سطح \ELim{S}
\item
  با استفاده از قضیه دیورژانس:
\end{itemize}

\begin{ELimportant}{}

\ELdisp{\oint_S\vect{E}\cdot d\vect{s}\ELbr =\frac{Q}{\epsilon_0}}

\end{ELimportant}

\end{ELjoin}

\begin{ELunit}

\begin{itemize}
\tightlist
\item
  این رابطه شکلی از \textbf{قانون گوس} است

  \begin{itemize}
  \tightlist
  \item
    شار کل خروجی شدت میدان الکتریکی از هر سطح بسته برابر با کل بار داخل سطح تقسیم بر \ELim{\epsilon_0} است.
  \end{itemize}
\item
  شکل انتگرالی معادلات دوم \ELim{\ELto} انتگرال‌گیری روی سطح باز دلخواه \ELim{S}
\end{itemize}

\ELdisp{\int_S(\nabla\times\vect{E})\cdot d\vect{s}\ELbr =0}

\end{ELunit}

\begin{ELjoin}

\begin{itemize}
\tightlist
\item
  با استفاده از قضیه استوکس:
\end{itemize}

\begin{ELimportant}{}

\ELdisp{\oint_C\vect{E}\cdot\dif\vect{\ell}\ELbr =0}

\end{ELimportant}

\end{ELjoin}

\begin{ELunit}

\begin{itemize}
\tightlist
\item
  انتگرال خطی عددی شدت میدان الکتریکی ساکن دور هر مسیر بسته صفر است.
\item
  حاصلضرب عددی \ELim{\vect{E}\cdot\dif\vect{\ell}} روی هر مسیری که انتگرال‌گیری شود، ولتاژ در امتداد آن مسیر است.
\item
  این معادله بیانی از \textbf{قانون ولتاژ کرشهف} است

  \begin{itemize}
  \tightlist
  \item
    مجموع جبری افت ولتاژهای اطراف هر مسیر بسته برابر با صفر است.
  \end{itemize}
\item
  انتگرال‌های خطی میدان غیرگردشی \ELim{\vect{E}} مستقل از مسیر بوده و فقط به نقاط ابتدایی و انتهایی بستگی دارد.
\end{itemize}

\ELfigure{m02-two-paths}{}

\ELdisp{\oint_C\vect{E}\cdot\dif\vect{\ell}\ELbr =0}
\ELdisp{\int_{C_1}\vect{E}\cdot\dif\vect{\ell}+\int_{C_2}\vect{E}\cdot\dif\vect{\ell}\ELbr =0}
\ELdisp{\int_{P_1}^{P_2}\vect{E}\cdot\dif\vect{\ell}+\int_{P_2}^{P_1}\vect{E}\cdot\dif\vect{\ell}\ELbr =0}
\ELdisp{\int_{\substack{P_1\\\text{\lr{Along }}C_1}}^{P_2}\vect{E}\cdot\dif\vect{\ell}\ELbr =-\int_{\substack{P_2\\\text{\lr{Along }}C_2}}^{P_1}\vect{E}\cdot\dif\vect{\ell}}
\ELdisp{\int_{\substack{P_1\\\text{\lr{Along }}C_1}}^{P_2}\vect{E}\cdot\dif\vect{\ell}\ELbr =\int_{\substack{P_1\\\text{\lr{Along }}C_2}}^{P_2}\vect{E}\cdot\dif\vect{\ell}}

\end{ELunit}

\ELhold

\ELsection{قانون کولمب}{قانون کولمب}

\begin{ELunit}

\ELfigure{m02-map-coulomb}{}

\begin{itemize}
\tightlist
\item
  \textbf{شدت میدان الکتریکی ناشی از بار نقطه‌ای \ELim{q} در فضای آزاد}

  \begin{itemize}
  \tightlist
  \item
    سطح کروی فرضی به شعاع \ELim{R} و مرکز \ELim{q} را در نظر می‌گیریم.
  \item
    چون بار نقطه‌ای دارای هیچ جهت خاصی نیست، انتظار داریم که میدان در همه جا شعاعی بوده و شدت آن برروی سطح کره ثابت باشد.
  \end{itemize}
\end{itemize}

\ELfigure{m02-point-charge}{}

\begin{itemize}
\tightlist
\item
  شدت میدان الکتریکی روی سطح کره:
\end{itemize}

\ELdisp{\vect{E}\ELbr =E_R\uvec{R}}
\ELdisp{\oint_S\vect{E}\cdot d\vect{s}\ELbr =\oint_S(\uvec{R}E_R)\cdot\uvec{R}\,ds\ELbr =\frac{q}{\epsilon_0}}
\ELdisp{E_R\oint_S ds\ELbr =E_R(4\pi R^2)\ELbr =\frac{q}{\epsilon_0}}

\end{ELunit}

\begin{ELimportant}{}

\ELdisp{\vect{E}\ELbr =\uvec{R}E_R\ELbr =\uvec{R}\frac{q}{4\pi\epsilon_0R^2}\qquad\ELbr (\mathrm{V/m})}

\end{ELimportant}

\begin{ELunit}

\begin{itemize}
\tightlist
\item
  اگر بار در مرکز دستگاه مختصات نباشد
\end{itemize}

\ELfigure{m02-offset-charge}{}

\ELdisp{\vect{E}_P\ELbr =\uvec{qP}\frac{q}{4\pi\epsilon_0\abs{\vect{R}-\vect{R}'}^2},\qquad\ELbr  \uvec{qP}\ELbr =\frac{\vect{R}-\vect{R}'}{\abs{\vect{R}-\vect{R}'}}}

\end{ELunit}

\begin{ELimportant}{}

\ELdisp{\vect{E}_P\ELbr =\frac{q(\vect{R}-\vect{R}')}{4\pi\epsilon_0\abs{\vect{R}-\vect{R}'}^3}\qquad\ELbr (\mathrm{V/m})}

\end{ELimportant}

\begin{ELjoin}

\begin{itemize}
\tightlist
\item
  هنگامیکه بار نقطه‌ای \ELim{q_2} در میدان الکتریکی بار نقطه‌ای دیگر \ELim{q_1} قرار می‌گیرد، نیروی وارد بر آن برابر است با
\end{itemize}

\begin{ELimportant}{بیان ریاضی قانون کولمب}

\ELdisp{\vect{F}_{12}\ELbr =q_2\vect{E}_{12}\ELbr =\uvec{R}\frac{q_1q_2}{4\pi\epsilon_0R^2}\qquad\ELbr (\mathrm{N})}

\end{ELimportant}

\end{ELjoin}

\begin{ELunit}

\ELfigure{m02-two-charges}{}

\end{ELunit}

\begin{ELjoin}

\begin{itemize}
\tightlist
\item
  \textbf{شدت میدان الکتریکی ناشی از مجموعه بارهای گسسته}
\end{itemize}

\begin{ELimportant}{}

\ELdisp{\vect{E}\ELbr =\frac{1}{4\pi\epsilon_0}\sum_{k=1}^{n}\frac{q_k(\vect{R}-\vect{R}'_k)}{\abs{\vect{R}-\vect{R}'_k}^3}\qquad\ELbr (\mathrm{V/m})}

\end{ELimportant}

\end{ELjoin}

\begin{ELunit}

\ELfigure{m02-discrete-charges}{}

\begin{itemize}
\tightlist
\item
  \textbf{شدت میدان الکتریکی ناشی از توزیع پیوسته بار}

  \begin{itemize}
  \tightlist
  \item
    میدان ناشی از بار دیفرانسیلی \ELim{dq} برابر است با
  \end{itemize}
\end{itemize}

\ELdisp{d\vect{E}\ELbr =\frac{dq(\vect{R}-\vect{R}')}{4\pi\epsilon_0\abs{\vect{R}-\vect{R}'}^3}}

\begin{itemize}
\tightlist
\item
  توزیع بار خطی
\end{itemize}

\ELdisp{dq\ELbr =\rho_\ell\,d\ell'\qquad\ELbr \qquad\ELbr  \vect{E}\ELbr =\int_{L'}\frac{\rho_\ell\,d\ell'(\vect{R}-\vect{R}')}{4\pi\epsilon_0\abs{\vect{R}-\vect{R}'}^3}}

\begin{itemize}
\tightlist
\item
  توزیع بار سطحی
\end{itemize}

\ELdisp{dq\ELbr =\rho_s\,ds'\qquad\ELbr \qquad\ELbr  \vect{E}\ELbr =\int_{S'}\frac{\rho_s\,ds'(\vect{R}-\vect{R}')}{4\pi\epsilon_0\abs{\vect{R}-\vect{R}'}^3}}

\begin{itemize}
\tightlist
\item
  توزیع بار حجمی
\end{itemize}

\ELdisp{dq\ELbr =\rho_v\,dv'\qquad\ELbr \qquad\ELbr  \vect{E}\ELbr =\int_{V'}\frac{\rho_v\,dv'(\vect{R}-\vect{R}')}{4\pi\epsilon_0\abs{\vect{R}-\vect{R}'}^3}}

\begin{itemize}
\tightlist
\item
  \textbf{مراحل محاسبه شدت میدان الکتریکی یک توزیع پیوسته بار}

  \begin{itemize}
  \tightlist
  \item
    تعیین دستگاه مختصات مناسب
  \item
    تعیین صحیح المان دیفرانسیلی \ELim{d\ell'}، \ELim{ds'} و \ELim{dv'} (دقت شود در این روابط اندازه بردار های دیفرانسیل طول و سطح استفاده شده است)
  \item
    تعیین صحیح بردار مکان نقطه مشاهده \ELim{(\vect{R})}
  \item
    تعیین صحیح بردار مکان نقطه منبع \ELim{(\vect{R}')}
  \item
    محاسبه انتگرال

    \begin{itemize}
    \tightlist
    \item
      دقت شود که نسبت به مختصات پریم‌دار انتگرال می‌گیریم.
    \end{itemize}
  \end{itemize}
\end{itemize}

\end{ELunit}

\ELnextnum{2-1}

\begin{ELexample}{}

شدت میدان الکتریکی یک بار خطی مستقیم و به طول بینهایت را با چگالی یکنواخت \ELim{\rho_\ell} تعیین کنید.

\ELfigure{m02-ex1}{}

\begin{ELsolution}

\ELdisp{\vect{E}\ELbr =\int_{L'}\frac{\rho_\ell\,d\ell'(\vect{R}-\vect{R}')}{4\pi\epsilon_0\abs{\vect{R}-\vect{R}'}^3}}

\begin{itemize}
\tightlist
\item
  انتخاب دستگاه مختصات استوانه‌ای
\end{itemize}

\ELdisp{d\ell'\ELbr =dz'\qquad\ELbr  \vect{R}\ELbr =r\uvec{r}\qquad\ELbr  \vect{R}'\ELbr =z'\uvec{z}}
\ELdisp{\vect{E}\ELbr =\frac{\rho_\ell}{4\pi\epsilon_0}\int_{-\infty}^{\infty}\frac{r\uvec{r}-z'\uvec{z}}{(r^2+z'^2)^{3/2}}\,dz'}

\begin{itemize}
\tightlist
\item
  به صورت شهودی واضح است که میدان مولفه در راستای \ELim{z} ندارد.
\end{itemize}

\ELdisp{\vect{E}\ELbr =\uvec{r}E_r\ELbr =\uvec{r}\frac{\rho_\ell r}{4\pi\epsilon_0}\int_{-\infty}^{\infty}\frac{dz'}{(r^2+z'^2)^{3/2}}}
\ELdisp{\vect{E}\ELbr =\uvec{r}\frac{\rho_\ell}{2\pi\epsilon_0r}\qquad\ELbr (\mathrm{V/m})}

\end{ELsolution}

\end{ELexample}

\ELnextnum{2-2}

\begin{ELexample}{}

حلقه‌ای به شعاع \ELim{a} و چگالی بار خطی \ELim{\rho_\ell} داریم. مطلوبست محاسبه شدت میدان الکتریکی روی محور حلقه به فاصله \ELim{h} از آن.

\ELfigure{m02-ex2}{}

\begin{ELsolution}

\ELdisp{\vect{E}\ELbr =\int_{L'}\frac{\rho_\ell\,d\ell'(\vect{R}-\vect{R}')}{4\pi\epsilon_0\abs{\vect{R}-\vect{R}'}^3}}

\begin{itemize}
\tightlist
\item
  انتخاب دستگاه مختصات استوانه‌ای
\end{itemize}

\ELdisp{d\ell'\ELbr =a\,d\phi'\qquad\ELbr  \vect{R}\ELbr =h\uvec{z}\qquad\ELbr  \vect{R}'\ELbr =a\uvec{r'}}
\ELdisp{\vect{E}\ELbr =\frac{\rho_\ell}{4\pi\epsilon_0}\int_0^{2\pi}\frac{h\uvec{z}-a\uvec{r'}}{(a^2+h^2)^{3/2}}\,a\,d\phi'}

\begin{itemize}
\tightlist
\item
  به صورت شهودی واضح است که میدان مولفه در راستای \ELim{r} ندارد.
\end{itemize}

\ELdisp{\vect{E}\ELbr =\frac{\rho_\ell ah}{4\pi\epsilon_0(a^2+h^2)^{3/2}}\uvec{z}\int_0^{2\pi}d\phi'\ELbr =\frac{\rho_\ell ah}{2\epsilon_0(a^2+h^2)^{3/2}}\uvec{z}}

\end{ELsolution}

\end{ELexample}

\ELnextnum{2-3}

\begin{ELexample}{}

شدت میدان الکتریکی یک بار مسطح بینهایت با چگالی بار سطحی یکنواخت \ELim{\rho_s} را تعیین کنید.

\ELfigure{m02-ex3}{}

\begin{ELsolution}

\ELdisp{\vect{E}\ELbr =\int_{S'}\frac{\rho_s\,ds'(\vect{R}-\vect{R}')}{4\pi\epsilon_0\abs{\vect{R}-\vect{R}'}^3}}

\begin{itemize}
\tightlist
\item
  انتخاب دستگاه مختصات استوانه‌ای
\end{itemize}

\ELdisp{ds'\ELbr =r'\,dr'\,d\phi'\qquad\ELbr  \vect{R}\ELbr =h\uvec{z}\qquad\ELbr  \vect{R}'\ELbr =r'\uvec{r'}}
\ELdisp{\vect{E}\ELbr =\frac{\rho_s}{4\pi\epsilon_0}\int_0^{2\pi}\!\!\int_0^{\infty}\frac{h\uvec{z}-r'\uvec{r'}}{(h^2+r'^2)^{3/2}}\,r'\,dr'\,d\phi'}

\begin{itemize}
\tightlist
\item
  به صورت شهودی واضح است که میدان مولفه در راستای \ELim{r} ندارد.
\end{itemize}

\ELdisp{\vect{E}\ELbr =\frac{\rho_sh\uvec{z}}{4\pi\epsilon_0}\int_0^{2\pi}\!\!\int_0^{\infty}\frac{r'\,dr'\,d\phi'}{(h^2+r'^2)^{3/2}}\ELbr =\frac{\rho_s}{2\epsilon_0}\uvec{z}}

\end{ELsolution}

\end{ELexample}

\ELhold

\ELsection{قانون گوس}{قانون گوس}

\begin{ELunit}

\ELfigure{m02-map-gauss}{}

\end{ELunit}

\begin{ELimportant}{}

\ELdisp{\oint_S\vect{E}\cdot d\vect{s}\ELbr =\frac{Q}{\epsilon_0}}

\end{ELimportant}

\begin{ELunit}

\begin{itemize}
\tightlist
\item
  \ELim{Q}: کل بار موجود در حجم \ELim{V} احاطه شده توسط سطح \ELim{S}
\item
  شار کل خروجی شدت میدان الکتریکی از هر سطح بسته برابر با کل بار داخل سطح تقسیم بر \ELim{\epsilon_0} است.
\item
  سطح \ELim{S} می‌تواند هر سطح بسته اختیاری باشد که به دلیل راحتی انتخاب شده است.
\end{itemize}

\end{ELunit}

\begin{ELremark}{}

قانون گوس همیشه برقرار است اما در استفاده از آن برای تعیین شدت میدان الکتریکی باید به نکته زیر دقت شود.

\begin{itemize}
\tightlist
\item
  اساس بکارگیری قانون گوس، برای تعیین شدت میدان الکتریکی، نخست در تشخیص شرایط تقارنی و دوم در انتخاب مناسب سطحی که روی آن مؤلفه عمودی \ELim{\vect{E}} ناشی از یک توزیع بار مشخص ثابت باشد (سطح گوسی)، است.
\end{itemize}

\end{ELremark}

\ELnextnum{2-4}

\begin{ELexample}{}

با استفاده از قانون گوس، شدت میدان الکتریکی ناشی از یک بار خطی مستقیم و بینهایت بلند را با چگالی یکنواخت \ELim{\rho_\ell} در فاصله \ELim{r} از آن تعیین کنید.

\ELfigure{m02-ex4}{}

\begin{ELsolution}

\ELdisp{\vect{E}\ELbr =\uvec{r}E_r}

\begin{itemize}
\tightlist
\item
  انتخاب یک استوانه به طول \ELim{L} و شعاع \ELim{r} به عنوان سطح گوسی

  \begin{itemize}
  \tightlist
  \item
    برای سطح بالایی و پایینی: \ELim{\vect{E}\cdot d\vect{s}=0}
  \item
    برای سطح جانبی: \ELim{d\vect{s}=\uvec{r}\,r\,d\phi\,dz}
  \end{itemize}
\end{itemize}

\ELdisp{\oint_S\vect{E}\cdot d\vect{s}\ELbr =\int_0^L\!\!\int_0^{2\pi}E_r\,r\,d\phi\,dz\ELbr =2\pi rLE_r}
\ELdisp{Q\ELbr =\rho_\ell L\qquad\ELbr  2\pi rLE_r\ELbr =\frac{\rho_\ell L}{\epsilon_0}\qquad\ELbr  \vect{E}\ELbr =\uvec{r}E_r\ELbr =\uvec{r}\frac{\rho_\ell}{2\pi\epsilon_0r}}

\end{ELsolution}

\end{ELexample}

\ELnextnum{2-5}

\begin{ELexample}{}

شدت میدان الکتریکی ناشی از یک بار مسطح بینهایت با چگالی بار سطحی یکنواخت \ELim{\rho_s} را تعیین کنید.

\ELfigure{m02-ex5}{}

\begin{ELsolution}

\ELdisp{\vect{E}\ELbr =\pm E_z\uvec{z}}

\begin{itemize}
\tightlist
\item
  انتخاب یک مکعب به عنوان سطح گوسی
\item
  برای سطح بالایی
\end{itemize}

\ELdisp{\vect{E}\cdot d\vect{s}\ELbr =(\uvec{z}E_z)\cdot(\uvec{z}\,ds)\ELbr =E_z\,ds}

\begin{itemize}
\tightlist
\item
  برای سطح پایینی
\end{itemize}

\ELdisp{\vect{E}\cdot d\vect{s}\ELbr =(-\uvec{z}E_z)\cdot(-\uvec{z}\,ds)\ELbr =E_z\,ds}

\begin{itemize}
\tightlist
\item
  برای سطوح جانبی
\end{itemize}

\ELdisp{\vect{E}\cdot d\vect{s}\ELbr =0}

\begin{itemize}
\tightlist
\item
  بنابراین
\end{itemize}

\ELdisp{\oint_S\vect{E}\cdot d\vect{s}\ELbr =2E_z\int_A ds\ELbr =2E_zA}
\ELdisp{Q\ELbr =\rho_sA\qquad\ELbr  2E_zA\ELbr =\frac{\rho_sA}{\epsilon_0}}

\begin{ELimportant}{}

\ELdisp{\vect{E}\ELbr =\uvec{z}E_z\ELbr =\uvec{z}\frac{\rho_s}{2\epsilon_0},\qquad\ELbr  z>0}
\ELdisp{\vect{E}\ELbr =-\uvec{z}E_z\ELbr =-\uvec{z}\frac{\rho_s}{2\epsilon_0},\qquad\ELbr  z<0}

\end{ELimportant}

\end{ELsolution}

\end{ELexample}

\ELnextnum{2-6}

\begin{ELexample}{}

میدان \ELim{\vect{E}} ناشی از یک ابر الکترونی با چگالی بار حجمی \ELim{\rho=-\rho_0} در ناحیه \ELim{0\leq R\leq b} و \ELim{\rho=0} در ناحیه \ELim{R>b} را تعیین کنید.

\ELfigure{m02-ex6}{}

\begin{ELsolution}

\ELdisp{\vect{E}\ELbr =\uvec{R}E_R}

\begin{itemize}
\tightlist
\item
  الف) \ELim{0\leq R\leq b}

  \begin{itemize}
  \tightlist
  \item
    انتخاب یک کره به عنوان سطح گوسی
  \end{itemize}
\end{itemize}

\ELdisp{d\vect{s}\ELbr =\uvec{R}\,ds}
\ELdisp{\oint_{S_i}\vect{E}\cdot d\vect{s}\ELbr =E_R\int_{S_i}ds\ELbr =E_R4\pi R^2}
\ELdisp{Q\ELbr =\int_V\rho\,dv\ELbr =-\rho_0\int_V dv\ELbr =-\rho_0\frac{4\pi}{3}R^3\qquad\ELbr  \vect{E}\ELbr =-\uvec{R}\frac{\rho_0}{3\epsilon_0}R}

\begin{itemize}
\tightlist
\item
  ب) \ELim{R\geq b}

  \begin{itemize}
  \tightlist
  \item
    انتخاب یک کره به عنوان سطح گوسی
  \end{itemize}
\end{itemize}

\ELdisp{d\vect{s}\ELbr =\uvec{R}\,ds}
\ELdisp{\oint_{S_o}\vect{E}\cdot d\vect{s}\ELbr =E_R\int_{S_o}ds\ELbr =E_R4\pi R^2}
\ELdisp{Q\ELbr =-\rho_0\frac{4\pi}{3}b^3\qquad\ELbr  \vect{E}\ELbr =-\uvec{R}\frac{\rho_0b^3}{3\epsilon_0R^2}}

\end{ELsolution}

\end{ELexample}

\ELnextnum{2-7}

\begin{ELexample}{}

یک توزیع بار کروی به شعاع \ELim{b} دارای چگالی زیر است. شدت میدان الکتریکی را در تمام نقاط فضا بدست آورید.
\ELdisp{\rho_v\ELbr =\begin{cases}\dfrac{\rho_0R}{b}&0\leq R\leq b\\\noalign{\vskip 2mm}0&R>b\end{cases}}

\ELfigure{m02-ex7}{}

\begin{ELsolution}

\ELdisp{\vect{E}\ELbr =\uvec{R}E_R}

\begin{itemize}
\tightlist
\item
  الف) \ELim{0\leq R\leq b}

  \begin{itemize}
  \tightlist
  \item
    انتخاب یک کره به عنوان سطح گوسی
  \end{itemize}
\end{itemize}

\ELdisp{d\vect{s}\ELbr =\uvec{R}\,ds}
\ELdisp{\oint_{S_i}\vect{E}\cdot d\vect{s}\ELbr =E_R\int_{S_i}ds\ELbr =E_R4\pi R^2}
\ELdisp{Q\ELbr =\int_v\rho_v\,dv\ELbr =\int_0^{2\pi}\!\!\int_0^{\pi}\!\!\int_0^{R}\frac{\rho_0R}{b}R^2\sin\theta\,dR\,d\theta\,d\phi\ELbr =\frac{\rho_0\pi R^4}{b}}
\ELdisp{\vect{E}\ELbr =E_R\uvec{R}\ELbr =\frac{\rho_0R^2}{4\epsilon_0b}\uvec{R}}

\begin{itemize}
\tightlist
\item
  ب) \ELim{R\geq b}

  \begin{itemize}
  \tightlist
  \item
    انتخاب یک کره به عنوان سطح گوسی
  \end{itemize}
\end{itemize}

\ELdisp{d\vect{s}\ELbr =\uvec{R}\,ds}
\ELdisp{\oint_{S_o}\vect{E}\cdot d\vect{s}\ELbr =E_R\int_{S_o}ds\ELbr =E_R4\pi R^2\qquad\ELbr  Q\ELbr =\rho_0\pi b^3}
\ELdisp{\vect{E}\ELbr =E_R\uvec{R}\ELbr =\frac{\rho_0b^3}{4\epsilon_0R^2}\uvec{R}}

\end{ELsolution}

\end{ELexample}

\ELhold

\ELsection{پتانسیل الکتریکی}{پتانسیل الکتریکی}

\begin{ELunit}

\ELfigure{m02-map-potential}{}

\begin{itemize}
\item
  مفهوم اختلاف پتانسیل الکتریکی بین دو نقطه و پتانسیل الکتریکی یک نقطه چیست؟
\item
  پتانسیل الکتریکی ناشی از یک بار نقطه‌ای، تعدادی بار نقطه‌ای گسسته و یک توزیع بار پیوسته چگونه محاسبه می‌شود؟
\item
  چگونه می‌توان با استفاده از پتانسیل الکتریکی، شدت میدان الکتریکی را محاسبه کرد و مزیت این کار نسبت به محاسبه مستقیم شدت میدان الکتریکی چیست؟
\item
  فرض کنید که می‌خواهیم بار نقطه‌ای \ELim{q} را که در میدان الکتریکی \ELim{\vect{E}} قرار دارد، به اندازه \ELim{d\ell} حرکت دهیم.
\end{itemize}

\ELfigure{m02-move-charge}{}

\begin{itemize}
\item
  نیروی وارد بر بار توسط میدان الکتریکی:
  \ELdisp{\vect{F}_E\ELbr =q\vect{E}}
\item
  اندازه مؤلفه این نیرو در راستای \ELim{d\ell}:
  \ELdisp{F_{EL}\ELbr =\vect{F}_E\cdot\uvec{\ell}\ELbr =q\vect{E}\cdot\uvec{\ell}}
\item
  نیرویی که باید به بار اعمال شود:
  \ELdisp{F_{apply}\ELbr =-F_{EL}\ELbr =-q\vect{E}\cdot\uvec{\ell}}
\item
  کار انجام شده برای حرکت بار به اندازه \ELim{d\ell}:
  \ELdisp{dW\ELbr =F_{apply}\,d\ell\ELbr =-q\vect{E}\cdot\dif\vect{\ell}}
\item
  کار انجام شده برای حرکت بار از نقطه \ELim{P_1} به نقطه \ELim{P_2}:
  \ELdisp{W\ELbr =-q\int_{P_1}^{P_2}\vect{E}\cdot\dif\vect{\ell}}
\end{itemize}

\ELdisp{\nabla\times\vect{E}\ELbr =0\ELbr \Rightarrow\vect{E}\ELbr =-\nabla V}
\ELdisp{\nabla V\cdot\uvec{\ell}\ELbr =\frac{dV}{d\ell}}
\ELdisp{q\ELbr =+1\,\mathrm{C}\qquad\ELbr  W\ELbr =\int_{P_1}^{P_2}\nabla V\cdot\dif\vect{\ell}\quad\ELbr \ELbr \Longrightarrow\quad\ELbr  W\ELbr =\int_{P_1}^{P_2}dV\ELbr =V_2-V_1}

\end{ELunit}

\begin{ELdefinition}{اختلاف پتانسیل الکتریکی}

اختلاف پتانسیل الکتریکی بین دو نقطه، کار انجام شده بوسیله منبع بیرونی در حرکت بار مثبت واحد از یک نقطه به نقطه دیگر است.
\ELdisp{V_2-V_1\ELbr =-\int_{P_1}^{P_2}\vect{E}\cdot\dif\vect{\ell}}

\begin{itemize}
\tightlist
\item
  به مسیر انتگرال‌گیری وابسته نیست
\end{itemize}

\end{ELdefinition}

\begin{ELunit}

\begin{itemize}
\tightlist
\item
  معمولاً در بسیاری از موارد یک نقطه مرجع در بینهایت با پتانسیل صفر در نظر می‌گیرند و اختلاف پتانسیل نقاط مختلف نسبت به این نقطه مرجع را بیان می‌کنند.
\end{itemize}

\end{ELunit}

\begin{ELimportant}{}

\ELdisp{V_P\ELbr =-\int_{\infty}^{P}\vect{E}\cdot\dif\vect{\ell}}

\end{ELimportant}

\begin{ELremark}{نکته ۱}

علامت منفی در رابطه پتانسیل برای مطابقت با این بحث که حرکت در خلاف جهت میدان باعث افزایش پتانسیل الکتریکی می‌شود لازم است.

\end{ELremark}

\begin{ELremark}{نکته ۲}

گرادیان \ELim{V} عمود بر سطوح \ELim{V} ثابت است و در نتیجه خطوط میدان الکتریکی عمود بر سطوح هم‌پتانسیل هستند.

\end{ELremark}

\begin{ELunit}

\begin{itemize}
\tightlist
\item
  \textbf{پتانسیل الکتریکی ناشی از بار نقطه‌ای \ELim{q} به فاصله \ELim{R} از آن}
\end{itemize}

\ELfigure{m02-potential-point}{}

\ELdisp{V\ELbr =-\int_{\infty}^{R}\left(\frac{q}{4\pi\epsilon_0R^2}\uvec{R}\right)\cdot dR\,\uvec{R}\ELbr =\frac{q}{4\pi\epsilon_0R}}

\begin{itemize}
\tightlist
\item
  \textbf{پتانسیل الکتریکی ناشی از \ELim{n} بار نقطه‌ای گسسته}
\end{itemize}

\ELdisp{V\ELbr =\frac{1}{4\pi\epsilon_0}\sum_{k=1}^{n}\frac{q_k}{\abs{\vect{R}-\vect{R}'_k}}}

\end{ELunit}

\begin{ELjoin}

\begin{itemize}
\item
  \ELim{\vect{R}}: بردار مکان نقطه مشاهده؛ \ELim{\vect{R}'_k}: بردار مکان نقاط منبع
\item
  \textbf{پتانسیل الکتریکی ناشی از توزیع پیوسته بار}

  \begin{itemize}
  \item
    توزیع بار خطی
    \ELdisp{V\ELbr =\frac{1}{4\pi\epsilon_0}\int_{L'}\frac{\rho_\ell\,d\ell'}{\abs{\vect{R}-\vect{R}'}}}
  \item
    توزیع بار سطحی
    \ELdisp{V\ELbr =\frac{1}{4\pi\epsilon_0}\int_{S'}\frac{\rho_s\,ds'}{\abs{\vect{R}-\vect{R}'}}}
  \item
    توزیع بار حجمی
    \ELdisp{V\ELbr =\frac{1}{4\pi\epsilon_0}\int_{V'}\frac{\rho_v\,dv'}{\abs{\vect{R}-\vect{R}'}}}
  \end{itemize}
\end{itemize}

\ELnextnum{2-8}

\begin{ELexample}{}

پتانسیل الکتریکی و شدت میدان الکتریکی ناشی از یک دوقطبی الکتریکی را در یک نقطه دلخواه به فاصله بسیار دور از آن بدست آورید.

\begin{ELsolution}

\ELdisp{V\ELbr =\frac{1}{4\pi\epsilon_0}\sum_{k=1}^{2}\frac{q_k}{\abs{\vect{R}-\vect{R}'_k}}\ELbr =\frac{1}{4\pi\epsilon_0}\left[\frac{q}{\abs{\vect{R}-\vect{d}/2}}+\frac{-q}{\abs{\vect{R}+\vect{d}/2}}\right]}

\ELfigure{m02-ex8-dipole}{}

\ELdisp{V\ELbr =\frac{q}{4\pi\epsilon_0}\left(\frac{1}{R_+}-\frac{1}{R_-}\right)}
\ELdisp{\frac{1}{R_+}\ELbr \cong\left(R-\frac{d}{2}\cos\theta\right)^{-1}\ELbr \cong R^{-1}\left(1+\frac{d}{2R}\cos\theta\right)}
\ELdisp{\frac{1}{R_-}\ELbr \cong\left(R+\frac{d}{2}\cos\theta\right)^{-1}\ELbr \cong R^{-1}\left(1-\frac{d}{2R}\cos\theta\right)}

\begin{itemize}
\tightlist
\item
  \ELim{\vect{p}=q\vect{d}}: بردار گشتاور دوقطبی الکتریکی
\end{itemize}

\ELdisp{V\ELbr =\frac{qd\cos\theta}{4\pi\epsilon_0R^2}\quad\ELbr \overset{\vect{d}=d\uvec{z}}{\Longrightarrow}\quad\ELbr  V\ELbr =\frac{q\vect{d}\cdot\uvec{R}}{4\pi\epsilon_0R^2}\quad\ELbr \overset{\vect{p}=q\vect{d}}{\Longrightarrow}\quad\ELbr  V\ELbr =\frac{\vect{p}\cdot\uvec{R}}{4\pi\epsilon_0R^2}}
\ELdisp{\begin{aligned}\vect{E}=-\nabla V&=-\uvec{R}\frac{\partial V}{\partial R}-\uvec{\theta}\frac{\partial V}{R\,\partial\theta}\\&=\frac{p}{4\pi\epsilon_0R^3}(\uvec{R}2\cos\theta+\uvec{\theta}\sin\theta)\end{aligned}}

\ELfigure{m02-ex8-field}{}

\end{ELsolution}

\end{ELexample}

\end{ELjoin}

\ELnextnum{2-9}

\begin{ELexample}{}

پتانسیل الکتریکی و شدت میدان الکتریکی ناشی از یک دیسک باردار به شعاع \ELim{b} و چگالی بار سطحی \ELim{\rho_s} را بر روی محور آن بدست آورید.

\ELfigure{m02-ex9}{}

\begin{ELsolution}

\ELdisp{V\ELbr =\frac{1}{4\pi\epsilon_0}\int_{S'}\frac{\rho_s\,ds'}{\abs{\vect{R}-\vect{R}'}}}
\ELdisp{ds'\ELbr =r'\,dr'\,d\phi'\qquad\ELbr  \vect{R}\ELbr =z\uvec{z}\qquad\ELbr  \vect{R}'\ELbr =r'\uvec{r'}}
\ELdisp{V\ELbr =\frac{\rho_s}{4\pi\epsilon_0}\int_0^{2\pi}\!\!\int_0^{b}\frac{r'\,dr'\,d\phi'}{\sqrt{z^2+r'^2}}\ELbr =\frac{\rho_s}{2\epsilon_0}\left[\sqrt{z^2+b^2}-z\right]\qquad\ELbr  z>0}
\ELdisp{\vect{E}\ELbr =-\nabla V\ELbr =\uvec{z}\frac{\rho_s}{2\epsilon_0}\left[1-\frac{z}{\sqrt{z^2+b^2}}\right]\qquad\ELbr  z>0}

\end{ELsolution}

\end{ELexample}

\ELnextnum{2-10}

\begin{ELexample}{}

پتانسیل الکتریکی و شدت میدان الکتریکی ناشی از یک بار خطی به طول \ELim{L} و چگالی بار \ELim{\rho_\ell} را در امتداد آن بدست آورید.

\ELfigure{m02-ex10}{}

\begin{ELsolution}

\ELdisp{V\ELbr =\frac{1}{4\pi\epsilon_0}\int_{L'}\frac{\rho_\ell\,d\ell'}{\abs{\vect{R}-\vect{R}'}}}
\ELdisp{d\ell'\ELbr =dz'\qquad\ELbr  \vect{R}\ELbr =z\uvec{z}\qquad\ELbr  \vect{R}'\ELbr =z'\uvec{z}}
\ELdisp{V\ELbr =\frac{\rho_\ell}{4\pi\epsilon_0}\int_{-L/2}^{L/2}\frac{dz'}{z-z'}\ELbr =\frac{\rho_\ell}{4\pi\epsilon_0}\ln\left[\frac{z+L/2}{z-L/2}\right]\qquad\ELbr  z>L/2}
\ELdisp{\vect{E}\ELbr =-\nabla V\ELbr =-\uvec{z}\frac{dV}{dz}\ELbr =\uvec{z}\frac{\rho_\ell L}{4\pi\epsilon_0\left[z^2-(L/2)^2\right]}\qquad\ELbr  z>L/2}

\end{ELsolution}

\end{ELexample}

\ELnextnum{2-11}

\begin{ELexample}{}

توزیع بار حجمی با چگالی \ELim{\rho=-\rho_0} در ناحیه \ELim{0\leq R\leq b} وجود دارد. پتانسیل الکتریکی در یک نقطه دلخواه درون این حجم چقدر است؟

\ELfigure{m02-ex11}{}

\begin{ELsolution}

\ELdisp{\vect{E}\ELbr =-\uvec{R}\frac{\rho_0}{3\epsilon_0}R,\qquad\ELbr  0\leq R\leq b}
\ELdisp{\vect{E}\ELbr =-\uvec{R}\frac{\rho_0b^3}{3\epsilon_0R^2},\qquad\ELbr  R\geq b}
\ELdisp{\begin{aligned}V=-\int_{\infty}^{R}\vect{E}\cdot\dif\vect{\ell}&=-\int_{\infty}^{b}\frac{-\rho_0b^3}{3\epsilon_0R^2}\uvec{R}\cdot dR\,\uvec{R}-\int_{b}^{R}\frac{-\rho_0R}{3\epsilon_0}\uvec{R}\cdot dR\,\uvec{R}\\&=-\frac{\rho_0b^2}{3\epsilon_0}-\frac{\rho_0}{6\epsilon_0}\left(b^2-R^2\right)\end{aligned}}

\end{ELsolution}

\end{ELexample}

\ELhold

\ELsection{رفتار میدان الکتریکی ساکن در محیط‌های مادی}{رفتار میدان الکتریکی ساکن در محیط‌های مادی}

\begin{ELunit}

\ELfigure{m02-map-media}{}

\begin{itemize}
\item
  تقسیم بندی مواد براساس خواص الکتریکی

  \begin{itemize}
  \tightlist
  \item
    هادی‌ها
  \item
    نیمه‌هادی‌ها
  \item
    عایق‌ها (دی‌الکتریک‌ها)
  \end{itemize}
\item
  در هادی‌ها الکترون‌های لایه‌های بیرونی اتم‌ها با نیروی ضعیفی نگه داشته شده‌اند و به سادگی از یک اتم به اتم دیگر منتقل می‌شوند.
\item
  اما الکترون‌ها در اتم عایق‌ها به شدت به مدار خود چسبیده‌اند و در شرایط عادی قادر به جدا شدن نمی‌باشند.
\item
  خواص الکتریکی نیمه‌هادی‌ها بین خواص الکتریکی هادی‌ها و عایق‌ها قرار دارد.
\item
  رفتار شدت میدان الکتریکی ساکن در حضور یک هادی چگونه است؟
\item
  رفتار شدت میدان الکتریکی ساکن در حضور یک دی الکتریک چگونه است؟
\item
  مفهوم چگالی شار الکتریکی چیست و چه رابطه ای با شدت میدان الکتریکی دارد؟
\item
  مفهوم ضریب دی الکتریک چیست؟
\item
  رفتار شدت میدان الکتریکی ساکن در مرز بین دو محیط متفاوت چگونه است؟
\end{itemize}

\end{ELunit}

\ELhold

\ELsection{هادی‌ها در میدان الکتریکی ساکن}{هادی‌ها در میدان الکتریکی ساکن}

\begin{ELunit}

\begin{itemize}
\tightlist
\item
  اگر بارهایی را درون یک هادی رها کنیم، میدانی در آن تشکیل می‌شود که باعث دور شدن بارها از هم و آمدن آن‌ها به سطح می‌شود.

  \begin{itemize}
  \item
    در شرایط سکون بار و میدان الکتریکی درون یک هادی صفر است.
    \ELdisp{\rho\ELbr =0\qquad\ELbr  \vect{E}\ELbr =0}
  \item
    زمان لازم برای توزیع بارها به سطح هادی و ایجاد حالت تعادل به میزان رسانندگی هادی بستگی دارد.
  \end{itemize}
\item
  اگر در شرایط سکون میدان الکتریکی روی سطح هادی دارای مؤلفه مماسی باشد، تولید یک نیروی مماسی نموده، بارها را به حرکت درآورده و حالت تعادل بار وجود نخواهد داشت.

  \begin{itemize}
  \tightlist
  \item
    بنابراین تحت شرایط سکون، میدان الکتریکی روی سطح یک هادی همه جا عمود بر سطح است. به عبارت دیگر تحت شرایط سکون سطح هادی یک سطح هم‌پتانسیل است. در واقع چون در تمام نقاط درون هادی \ELim{\vect{E}=0} است، کل هادی دارای پتانسیل الکتریکی یکسانی است.
  \end{itemize}
\item
  \textbf{شرایط مرزی در مرز بین هادی/فضای آزاد}
\end{itemize}

\ELfigure{m02-conductor-boundary}{}

\end{ELunit}

\begin{ELjoin}

\begin{itemize}
\item
  مؤلفه عمودی:
  \ELdisp{\oint_S\vect{E}\cdot d\vect{s}\ELbr =E_n\Delta S\ELbr =\frac{\rho_s\Delta S}{\epsilon_0}\qquad\ELbr \ELbr \Longrightarrow\qquad\ELbr  E_n\ELbr =\frac{\rho_s}{\epsilon_0}}
\item
  مؤلفه مماسی:
  \ELdisp{\oint\vect{E}\cdot d\vect{L}\ELbr =0}
  \ELdisp{E_t\Delta w-E_{N,\text{\lr{at }}b}\tfrac12\Delta h+E_{N,\text{\lr{at }}a}\tfrac12\Delta h\ELbr =0\qquad\ELbr \ELbr \Longrightarrow\qquad\ELbr  E_t\ELbr =0}
\end{itemize}

\ELnextnum{2-12}

\begin{ELexample}{}

بار نقطه‌ای \ELim{+Q} در مرکز یک پوسته کروی هادی به شعاع داخلی \ELim{R_i} و شعاع خارجی \ELim{R_o} قرار دارد. \ELim{\vect{E}} و \ELim{V} را در کل فضا بدست آورید.

\ELfigure{m02-ex12}{}

\begin{ELsolution}

\ELdisp{\vect{E}\ELbr =E_R\uvec{R}}

\begin{itemize}
\tightlist
\item
  به ازای \ELim{R>R_o}
\end{itemize}

\ELdisp{\oint_S\vect{E}\cdot d\vect{s}\ELbr =E_{R1}4\pi R^2\ELbr =\frac{Q}{\epsilon_0}}
\ELdisp{E_{R1}\ELbr =\frac{Q}{4\pi\epsilon_0R^2}\qquad\ELbr  V_1\ELbr =-\int_{\infty}^{R}(E_{R1})\,dR\ELbr =\frac{Q}{4\pi\epsilon_0R}}

\begin{itemize}
\tightlist
\item
  به ازای \ELim{R_i<R<R_o}
\end{itemize}

\ELdisp{E_{R2}\ELbr =0\qquad\ELbr  V_2\ELbr =V_1\Big|_{R=R_o}\ELbr =\frac{Q}{4\pi\epsilon_0R_o}}

\begin{itemize}
\tightlist
\item
  به ازای \ELim{R<R_i}
\end{itemize}

\ELdisp{E_{R3}\ELbr =\frac{Q}{4\pi\epsilon_0R^2}}
\ELdisp{V_3\ELbr =-\int_{\infty}^{R_o}E_{R1}\,dR-\int_{R_o}^{R_i}E_{R2}\,dR-\int_{R_i}^{R}E_{R3}\,dR}
\ELdisp{V_3\ELbr =\frac{Q}{4\pi\epsilon_0}\left(\frac1R+\frac{1}{R_o}-\frac{1}{R_i}\right)}

\ELfigure{m02-ex12-plots}{}

\end{ELsolution}

\end{ELexample}

\end{ELjoin}

\ELhold

\ELsection{دی‌الکتریک‌ها در میدان الکتریکی ساکن}{دی‌الکتریک‌ها در میدان الکتریکی ساکن}

\begin{ELunit}

\begin{itemize}
\tightlist
\item
  دی‌الکتریک‌های ایده‌آل دارای بار آزاد نیستند. در نتیجه مانند هادی‌ها چگالی بار و میدان الکتریکی داخلی آن‌ها برابر با صفر نیست.
\item
  دی‌الکتریک‌ها دارای \textbf{بارهای مقید} بوده و اعمال میدان خارجی به دی‌الکتریک باعث ایجاد دو قطبی‌های الکتریکی شده و ماده دی‌الکتریک را قطبی می‌کند.
\item
  دو قطبی‌های الکتریکی القا شده میدان الکتریکی را در داخل و خارج ماده دی‌الکتریک تغییر می‌دهند.
\end{itemize}

\ELfigure{m02-polarization}{}

\begin{itemize}
\tightlist
\item
  \textbf{توزیع‌های بار معادل در دی‌الکتریک‌های قطبی شده}

  \begin{itemize}
  \tightlist
  \item
    برای تجزیه و تحلیل تأثیر ماکروسکوپی دوقطبی‌های القا شده، بردار قطبی‌شدگی \ELim{\vect{P}} را به صورت زیر تعریف می‌کنیم
  \end{itemize}
\end{itemize}

\end{ELunit}

\begin{ELdefinition}{}

\ELdisp{\vect{P}\ELbr =\lim_{\Delta v\to0}\frac{\sum_{k=1}^{n\Delta v}\vect{p}_k}{\Delta v}\qquad\ELbr (\mathrm{C/m^2})}

\begin{itemize}
\tightlist
\item
  \ELim{n}: تعداد دوقطبی ها در واحد حجم؛ \ELim{\vect{p}_k}: گشتاور دوقطبی
\end{itemize}

\end{ELdefinition}

\begin{ELunit}

\begin{itemize}
\tightlist
\item
  اگر \ELim{d\vect{p}} گشتاور دو قطبی در یک حجم کوچک \ELim{dv'} باشد، آنگاه
\end{itemize}

\ELdisp{d\vect{p}\ELbr =\vect{P}\,dv'\qquad\ELbr  dV\ELbr =\frac{\vect{P}\cdot\uvec{R}}{4\pi\epsilon_0R^2}\,dv'\qquad\ELbr  V\ELbr =\frac{1}{4\pi\epsilon_0}\int_{V'}\frac{\vect{P}\cdot\uvec{R}}{R^2}\,dv'}
\ELdisp{V\ELbr =\frac{1}{4\pi\epsilon_0}\oint_{S'}\frac{\vect{P}\cdot\vect{a}'_n}{R}\,ds'+\frac{1}{4\pi\epsilon_0}\int_{V'}\frac{(-\nabla'\cdot\vect{P})}{R}\,dv'}

\begin{itemize}
\tightlist
\item
  پتانسیل الکتریکی و بنابراین شدت میدان الکتریکی ناشی از دی‌الکتریک قطبی شده می‌تواند از اثر توزیع‌های دو بار سطحی و حجمی به ترتیب با چگالی‌های زیر که آن‌ها را \textbf{چگالی بار قطبی‌شدگی} یا \textbf{چگالی بار مقید} می‌نامند، محاسبه شود
\end{itemize}

\end{ELunit}

\begin{ELimportant}{}

\ELdisp{\rho_{ps}\ELbr =\vect{P}\cdot\uvec{n}\qquad\ELbr \qquad\ELbr  \rho_p\ELbr =-\nabla\cdot\vect{P}}
\ELdisp{V\ELbr =\frac{1}{4\pi\epsilon_0}\oint_{S'}\frac{\rho_{ps}}{R}\,ds'+\frac{1}{4\pi\epsilon_0}\int_{V'}\frac{\rho_p}{R}\,dv'}

\end{ELimportant}

\begin{ELremark}{}

چون با یک جسم دی‌الکتریک خنثی از نظر الکتریکی سروکار داریم، بار کل جسم پس از قطبی شدن باید همچنان صفر باشد
\ELdisp{\begin{aligned}\text{\lr{Total charge}}&=\oint_S\rho_{ps}\,ds+\int_V\rho_p\,dv\\&=\oint_S\vect{P}\cdot\uvec{n}\,ds-\int_V\nabla\cdot\vect{P}\,dv=0\end{aligned}}

\end{ELremark}
